\scanpage{59}
où $u_{\unsure{E|\varepsilon}}$ est la \deleted{\unsure{autom}} perm. non identique de $E$ (\deleted{\unsure{toutes}} les fibres sont canoniquemt isom par les itérés de $u$)

b) Il y a une seule orbite dans $B$. Cette situation est aussi uniquement déterminée (modulo \inserted{$I$-}isom \deleted{\unsure{x}} une fois $I,u_I$ donnés, modulo isom une fois $n=\operatorname{card}I$ donné).

% slanted
\margin{``cas tordu''}
\diagram{\(\begin{array}{c@{\;}cccc}
 & &ub_1&&u^{\inserted{n-1}}\deleted{\unsure{u}}b_n\\[-2pt]
 & &\diagup\!\diagup&&\diagup\!\diagup\\[-2pt]
 &b_1&b_2&\ldots&b_n\\
 ub_n\,{=}&b'_1&b'_2&\ldots&b'_n\\
 \cline{2-5}
 &i_1&i_2&\ldots&i_n\\[-3pt]
 &&\shortparallel&&\shortparallel\\[-2pt]
 &&ui_1&&u^{\deleted{\unsure{n}}n-1}i_1
\end{array}\)}
\margin{\uline{NB} Classification des flèches $B\to\deleted{\unsure{B|\overline I}}I$ (\deleted{\unsure{surj.|$\operatorname{card}I$}}) d'\deleted{ens} $\mathbb Z$-ens. à \unsure{esp.|sp.} finis}

Ainsi, pour $n=\operatorname{Card}I$ fixé, les classes d'isom d'autom d'un $n$-cube (i.e. les classes de conjugaison d'éléments de $\operatorname{Aut}(C_n)=\{\pm1\}^n\cdot\unsure{\mathfrak S_n|\mathscr S_n}$) sont donnés par 2 suites d'entiers naturels
\[
\begin{aligned}
\alpha_1^+\ldots\alpha_n^+&\in\mathbb N\\
\alpha_1^-\ldots\alpha_n^-&\in\mathbb N
\end{aligned}
\]
\margin{$\alpha_i^+$ \inserted{($\alpha_i^-$) $=$} nb de cycles de long. $i$ de $u_I$ au dessus duquel il y a \unsure{un cycle} de long $i$ (resp. $2i$)}
satisfaisant
\[
\sum_1^{\unsure{n|\relbar}}i(\underbrace{\alpha_i^++\alpha_i^-}_{\alpha_i})=n
\]

\uline{Étude des facettes fixes} $(\ne\varnothing)$ \deleted{\illegible{1}} Correspondent aux parties non vides \inserted{$B'$} de $B$ qui a) stables par $I$ b) Sections partielles. Au dessus de chaque $I_\alpha$ correspondant au ``\unsure{bicycle|cycle|cas tordu}'',
% archive sheet number: 58
